\chapter*{Avant-propos}
\markboth{AVANT-PROPOS}{}
\addcontentsline{toc}{chapter}{Avant-propos}

\section*{\`A qui s'adresse ce guide}

Ce guide s'adresse aux trois profils qui exploitent la plateforme au
quotidien :

\begin{itemize}[leftmargin=*]
  \item l'\textbf{analyste}, qui suit les interventions, tranche sur les gestes
        que la plateforme refuse d'engager seule et annule ce qui doit l'\^etre ;
  \item l'\textbf{administrateur}, qui installe, configure, r\'edige les
        consignes de r\'eponse, g\`ere les comptes et r\'earme le coupe-circuit ;
  \item le \textbf{d\'ecideur}, qui consulte le portefeuille d'incidents et
        \'edite les rapports destin\'es \`a la hi\'erarchie ou \`a une
        proc\'edure.
\end{itemize}

Aucune connaissance de la programmation n'est requise pour les chapitres~1
\`a~7. Le chapitre~8, consacr\'e aux technologies employ\'ees, s'adresse
davantage \`a un profil technique ou \`a un \'evaluateur.

\section*{Ce que la plateforme fait}

CIRTDEFENSE re\c{c}oit les observations des capteurs de s\'ecurit\'e du parc
sur\-veill\'e, les rapproche du catalogue des menaces du Centre, d\'ecide de la
conduite \`a tenir et \textbf{applique elle-m\^eme les gestes de confinement},
sans attendre l'accord pr\'ealable d'un agent. Elle informe ensuite l'\'equipe
d'astreinte, mesure l'effet de ce qu'elle a fait, et annule d'elle-m\^eme un
confinement qui aurait d\'egrad\'e le service qu'il prot\'egeait.

Cette posture --- agir d'abord, rendre compte ensuite --- est le choix
structurant du projet. Elle a une contrepartie technique strictement
appliqu\'ee : \textbf{la plateforme n'engage seule que ce qu'elle sait
annuler enti\`erement}. Tout geste \`a effet durable lui est interdit sans
d\'ecision humaine.

\section*{Ce que la plateforme ne fait pas}

Il est aussi utile de savoir o\`u elle s'arr\^ete :

\begin{itemize}[leftmargin=*]
  \item elle \textbf{n'agit pas} sur une menace qu'aucun document de
        r\'ef\'erence ne d\'ecrit --- elle le dit et s'abstient ;
  \item elle \textbf{n'ex\'ecute jamais} un geste irr\'eversible de sa propre
        initiative ;
  \item elle \textbf{ne remplace pas} l'analyste : elle lui \'evite les gestes
        d'urgence r\'ep\'etitifs et lui rend son temps pour l'investigation ;
  \item elle \textbf{ne d\'ecide rien au moyen d'un mod\`ele de langage}. Ce
        point est d\'etaill\'e au chapitre~8.
\end{itemize}

\section*{Comment lire ce guide}

Les chapitres~2 et~3 se lisent une fois, \`a l'installation. Les chapitres~4
\`a~7 sont des chapitres de r\'ef\'erence : on y revient au besoin. Le
chapitre~5, qui d\'eroule les sc\'enarios d'usage courants, est celui par lequel
commencer si vous d\'ecouvrez la plateforme d\'ej\`a install\'ee.
